앞 장에서 그래프의 정의와 저장 방법, 그리고 용어들을 알아보았다. 이번 장에서는 그래프의 큰 분류와 특수한 그래프 형태 몇 가지를 소개하고, 각각이 어떻게 이용되는지 설명한다.

\section{그래프의 분류}

\subsection{방향성}
일반적으로 이전까지 살펴보았던 그래프의 간선은 두 정점의 연결 관계를 나타내는 것으로, 양방향 이동이 가능하다고 가정하였다. 이러한 그래프를 무방향 그래프라 부르고, 무방향 그래프의 간선은 무방향 간선이라고 한다.
하지만 그래프의 간선이 두 정점의 연결 관계가 아닌 이동 관계를 나타낼 수도 있다. 이러한 경우에는 $v$와 $u$를 연결하는 간선이 존재하더라도 $V$에서 $u$로의 이동은 허용되지만, $u$에서 $v$로의 이동은 허용되지 않을 수 있다.
이러한 그래프를 방향 그래프라고 하며, 방향 그래프의 간선을 방향 간선이라고 한다.

많은 문제에서 무방향 그래프를 다루게 될 것이므로 앞으로의 그래프 알고리즘 설명에서 무방향 그래프를 가정한다. 또한, 일반적으로 무방향 그래프에서 구현된 알고리즘을 방향 그래프로 바꾸는 것이 어렵지 않으므로 방향 그래프에 대한 논의는 생략하겠다(위상 정렬에서만 특수히 방향 그래프를 다룰 것이다).

\subsection{가중치}
일반적인 그래프에는 가중치가 존재하지 않다. 하지만 떄떄로 간선에 가중치가 존재하는 경우가 있으며, 이 가중치는 최단경로에만 영향을 준다. 앞에서 최단경로를 경로 중에 가장 적은 개수의 간선으로 이루어진 것으로 정의했지만, 가중치가 있는 경우에는 경로를 이루는 간선의 가중치의 합이 최소인 것이다.
가중치가 포함될 경우에는 최단 거리 이외의 성질은 달라지지 않으므로 가중치가 포함된 그래프는 최단거리 알고리즘 및 최소 스패닝 트리에서만 언급될 것이다.

\section{특수한 그래프}

이전 장에서 컴포넌트의 개념을 배웠다. 임의의 두 정점이 연결된 $V$의 부분집합을 컴포넌트라고 하는데, $V$ 자체가 컴포넌트인 그래프를 연결 그래프라고 한다. 다시 말해 전체 그래프의 임의의 두 정점 사이가 연결되어 있으면, 그 그래프는 연결 그래프라고 부른다.

\begin{definition}[연결 그래프]
    $G=(V, E)$에서 $V$가 컴포넌트이면 $G$를 \textbf{연결 그래프}라고 한다.
\end{definition}

정점 $N$개로 이루어진 무방향 연결 그래프 중 가능한 간선 개수 $M$의 최솟값은 얼마일까? 바로 $N-1$이다. 적어도 $N$개의 정점이 연결되려면 $N-1$개의 간선이 필요하다. 트리가 바로 이러한 그래프이다. 다시 말해 정점 $N$개와 간선 $N-1$개로 이루어진 연결 그래프를 트리라고 한다.

\begin{definition}[트리]
    $G=(V, E)$에서 $V$의 크기가 $E$의 크기에 1을 더한 값이면 $G$를 \textbf{트리}라고 한다.
\end{definition}

왜 트리라고 부를까? 트리는 하나의 정점에서 시작하여 위로 뻗어나가는 나무 모양의 구조를 이루기 때문이다. 이러한 의미에서 트리라는 이름이 붙었으며, 계층 구조나 상하 구조와 같이 나뭇가지처럼 뻗어나가야 하는 경우 트리를 통해 나타내는 경우가 많다.

트리에는 매우 많은 성질이 있다. 특히, 어떤 그래프 $G$가 트리인 것과 다음 명제들은 필요충분조건이다.
\begin{itemize}
    \item $G$가 무방향 연결 그래프이고, 간선 수에 1을 더하면 정점 수가 된다.
    \item $G$가 무방향 연결 그래프이고, 두 정점 사이를 잇는 경로는 유일하다.
    \item $G$가 무방향 연결 그래프이고, 사이클이 없다.
\end{itemize}

임의의 간선이 둘 이상의 사이클에 포함되지 않는 그래프를 생각해보자. 이러한 그래프는 트리는 아니지만 뾰족뾰족한 선인장과 같은 모양을 갖는다. 이러한 그래프를 선인장 그래프라고 한다.

\begin{definition}[선인장 그래프]
    $G=(V, E)$에서 $E$의 임의의 원소(간선)이 많아야 하나의 사이클에 속하면 $G$를 \textbf{선인장 그래프}라고 한다.
\end{definition}

선인장 그래프의 판별은 중요한 문제 중 하나이다. 이러한 문제는 이후에 이중 연결 요소와 같이 살펴볼 것이다.

중복 간선이 없어야 한다는 가정이 있다면, 무방향 연결 그래프에 존재할 수 있는 간선의 최대 개수는 몇 개일까? 당연히 모든 정점쌍이 간선으로 연결된 경우에 $M=N(N+1)/2$까지 가능할 것이다. 이러한 그래프를 완전 그래프라고 한다.

\begin{definition}[완전 그래프]
    $G=(V, E)$에서 $V$에 속하는 임의의 정점 $u, v$에 대해 $(u, v)$ 간선이 $E$에 존재하면 $G$를 \textbf{완전 그래프}라고 한다.
\end{definition}

때떄로 어떤 그래프는 양쪽 집합 사이를 연결하기도 한다. 이 경우에는 각 집합 내부의 연결 간선이 없을 것이다. 이러한 그래프를 이분 그래프라고 한다.

\begin{definition}[이분 그래프]
    어떤 그래프 $G=(V, E)$에서 $V$를 정확히 두 집합 $A$와 $B$로 $A \subset G, B \subset G, A \cup B=G, A \cap B = \phi$가 되도록 나눌 때, $A$의 임의의 두 원소 $u, v$에 대해 $(u, v) \notin E$가 성립하고 $B$의 임의의 두 원소 $u, v$에 대해 $(u, v) \notin E$가 성립하는 $A$와 $B$가 존재할 때 $G$를 \textbf{이분 그래프}라고 한다.
\end{definition}

우리는 이분 그래프에서 두 집합의 원소를 간선으로 연결된 경우에 한해 매칭할 수 있다(당연히 같은 원소는 두 번 이상 매칭되면 안될 것이다). 이러한 문제를 이분 매칭 문제라고 하며 이분 매칭 문제는 유량 문제의 시작점이 되는 문제로써 중요하다.

정의역과 공역이 같은 함수 $f$를 그래프로 나타낼 수 있다. 정의역과 공역을 정점 집합 $V$로 놓고, 각 $V$의 원소 $v$에 대해 $v$에서 $f(v)$로 가는 방향 간선을 설계하면 된다. 역으로 말하면 각 정점 $v$에 대해 진출 차수가 1인 모든 그래프는 적절한 함수 $f$를 찾을 수 있으므로 함수형 그래프라고 말할 수 있다.

\begin{definition}[함수형 그래프]
    $G=(V, E)$가 방향 그래프이고 임의의 정점 $v \in V$에 대해 진출 차수가 1이면 $G$를 \textbf{함수형 그래프}라고 한다.
\end{definition}

함수형 그래프는 $f$의 합성함수를 처리하는데 유용하다. 단순히 $f(f(f(v)))$를 구하고 싶다면 $v$에서 간선을 따라 세 번 이동하면 되기 때문이다. 이러한 함수형 그래프는 희소 배열을 시각적으로 나타낸 것과 같다.

다른 종류의 특수한 방향 그래프도 있다. 만약, 어떤 방향 그래프가 사이클을 갖지 않는다면, 이러한 그래프를 방향 비순환 그래프(DAG)라고 한다. 

\begin{definition}[방향 비순환 그래프]
    $G=(V, E)$가 방향 그래프이고, 이 그래프에 어떠한 사이클도 존재하지 않는다면 $G$를 \textbf{방향 비순환 그래프}라고 한다.
\end{definition}

어떤 그래프가 DAG인 것과 위상 정렬이 존재하는 것은 필요충분조건이다. 이와 관련된 자세한 내용은 위상 정렬 단원에서 살펴보겠다.

어떤 그래프는 간선끼리 교점이 존재하지 않도록 평면에 나타낼 수 있다. 즉, 주어진 그래프를 평면의 분할로 나타낼 수 있는 그래프가 존재한다. 이러한 그래프를 평면 그래프라고 한다.

\begin{definition}[평면 그래프]
    $G=(V, E)$가 간선끼리 교점이 존재하지 않도록 한 평면에 정점과 그 정점을 잇는 간선을 그릴 수 있을 때 $G$를 \textbf{평면 그래프}라고 한다. 매우 엄밀하게 말하면 평면에 graph embedding이 가능한 그래프이다.
\end{definition}

당연하게도 평면 그래프는 오일러의 등식 $v - e + f = 2$을 만족시킨다.

\section{연습문제 A}

\begin{problem}
    \Cref{code: adj-matrix}와 \Cref{code: adj-list}를 참고하여 방향 그래프의 인접행렬과 인접리스트를 구하는 코드를 작성하라.
\end{problem}
\begin{problem}
    \Cref{code: adj-matrix}와 \Cref{code: adj-list}를 참고하여 가중치가 있는 무방향 그래프의 인접행렬과 인접리스트를 구하는 코드를 작성하라.
\end{problem}
\begin{problem}
    정점 $N$개로 이루어진 무방향 연결 그래프의 간선 개수 최솟값이 $N-1$임을 증명하라. (hint. 수학적 귀납법)
\end{problem}
\begin{problem}
    트리의 성질로 주어진 세 명제가 동치임을 증명하라.
\end{problem}
\begin{problem}
    정점의 개수와 간선의 개수가 같은 연결 그래프(선인장)에 정확히 하나의 사이클이 존재하고, 그 사이클에 연결된 부분 그래프는 트리임을 증명하라.
\end{problem}
\begin{problem}
    임의의 연결 그래프에서 간선 몇 개를 적절히 제거하여 트리를 만들 수 있음을 보여라(이러한 트리를 스패닝 트리라고 한다).
\end{problem}
\begin{problem}
    graph embedding의 정의를 명확히 찾아보고 평면 그래프와 어떠한 관련이 있는지 분석하라.
\end{problem}
\begin{problem}
    평면 그래프의 필요충분조건을 말해주는 \textbf{쿠라토프스키 정리}가 있다. 이 정리를 이해하고 증명을 읽어보자.
\end{problem}

\section{연습문제 B}

\subsection{기초문제}
\begin{problem} \url{https://www.acmicpc.net/problem/9644} \end{problem}
\begin{problem} \url{https://www.acmicpc.net/problem/15166} \end{problem}
\begin{problem} \url{https://www.acmicpc.net/problem/22100} \end{problem}
\begin{problem} \url{https://www.acmicpc.net/problem/22392} \end{problem}
\begin{problem} \url{https://www.acmicpc.net/problem/23039} \end{problem}
\begin{problem} \url{https://www.acmicpc.net/problem/30411} \end{problem}

\subsection{응용문제}
\begin{problem} \url{https://www.acmicpc.net/problem/3461} \end{problem}
\begin{problem} \url{https://www.acmicpc.net/problem/4634} \end{problem}
\begin{problem} \url{https://www.acmicpc.net/problem/9666} \end{problem}
\begin{problem} \url{https://www.acmicpc.net/problem/13331} \end{problem}
\begin{problem} \url{https://www.acmicpc.net/problem/15010} \end{problem}
\begin{problem} \url{https://www.acmicpc.net/problem/21002} \end{problem}
\begin{problem} \url{https://www.acmicpc.net/problem/23154} \end{problem}
\begin{problem} \url{https://www.acmicpc.net/problem/24689} \end{problem}
\begin{problem} \url{https://www.acmicpc.net/problem/28111} \end{problem}

\subsection{도전문제}
\begin{problem} \url{https://www.acmicpc.net/problem/10591} \end{problem}
\begin{problem} \url{https://www.acmicpc.net/problem/12625} \end{problem}
\begin{problem} \url{https://www.acmicpc.net/problem/14519} \end{problem}
\begin{problem} \url{https://www.acmicpc.net/problem/16746} \end{problem}
\begin{problem} \url{https://www.acmicpc.net/problem/17772} \end{problem}
\begin{problem} \url{https://www.acmicpc.net/problem/18447} \end{problem}
\begin{problem} \url{https://www.acmicpc.net/problem/19220} \end{problem}
\begin{problem} \url{https://www.acmicpc.net/problem/19445} \end{problem}
\begin{problem} \url{https://www.acmicpc.net/problem/19552} \end{problem}
\begin{problem} \url{https://www.acmicpc.net/problem/20573} \end{problem}
\begin{problem} \url{https://www.acmicpc.net/problem/20966} \end{problem}
\begin{problem} \url{https://www.acmicpc.net/problem/24133} \end{problem}   
\begin{problem} \url{https://www.acmicpc.net/problem/25355} \end{problem}
\begin{problem} \url{https://www.acmicpc.net/problem/28332} \end{problem}